\section{Exact computation, additional terms, and reproducibility}\label{sec:computation}
The all-order assertions of this article are proved algebraically or
analytically; a finite computation is not their substitute. Computation is
nevertheless useful for certifying transcription, signs, bracket
conventions, and the additional explicit terms, and for detecting errors in
the finite displays. This section describes the exact computations that
accompany the article.

\subsection{Representation and independently organized calculations}
An associative polynomial is stored as a dictionary from words in $\{X,Y\}$
to exact rational numbers (Python's \texttt{fractions.Fraction}); the empty
word represents $1$. Multiplication concatenates words and discards
degrees beyond a requested cutoff. Exponential and logarithm are evaluated
by their defining series, which become finite under this truncation.
Because no relation is imposed on the two generators, equality of the
resulting dictionaries is equality in the truncated free associative
algebra, not a test on selected numerical matrices in which independent
commutators might collapse.

The three source treatments each supplied a standard-library Python
verifier, and all three are retained unchanged in the repository
(\texttt{paper-1/code}, \texttt{paper-2/code}, \texttt{paper-3/code}). Between
them they perform the following independently organized calculations, each
compared coefficient by coefficient with the reference expansion
$\log(e^Xe^Y)$:
\begin{itemize}
\item the finite word-cut formula \eqref{eq:wordcut} evaluated by a dynamic
program over all contiguous cuts of every word, including words with zero
coefficient (through degree $10$; $2046$ words);
\item the Poincar\'e operator-block expansion \eqref{eq:operatorblock},
assembled by repeated application of $e^{\ad_X}e^{t\ad_Y}-I$ with exact
integration of powers of $t$ (through degree $10$);
\item the right Dynkin projection \eqref{eq:dynkinword}, verified to fix
every homogeneous component (through degree $12$);
\item the total-degree Bernoulli recursion \eqref{eq:homrec} (through
degree $10$), the single-$Y$ Bernoulli series \eqref{eq:linearY} (through
degree $12$), and the quadratic kernel \eqref{eq:H2kernel} (through degree
$10$);
\item the primitive coproduct condition, by explicit computation of the
reduced unshuffle coproduct (through degree $10$);
\item the exchange parity \eqref{eq:parity} and word-reversal parity
\eqref{eq:reversal} (through degree $12$), and three-letter associativity
\eqref{eq:assoc} (through degree $6$);
\item the displayed terms \eqref{eq:z1}--\eqref{eq:z6}, entered
independently as nested brackets and expanded by \eqref{eq:bracketexpand};
\item the Zassenhaus exponents computed three ways, by residual extraction
\eqref{eq:residual}, by the Lie recurrence \eqref{eq:frec}, and by the
forward partition recurrence \eqref{eq:Cpartitions}, compared with each
other and with the explicit enumeration of the weighted-tree formula
\eqref{eq:treeformula} (through degree $10$), followed by reconstruction of
the ordered product and comparison with $e^{t(X+Y)}$ (through degree
$12$), and by the displayed $C_2,C_3,C_4$;
\item the ordered-simplex expansion \eqref{eq:simplex} (through degree
$10$), the symmetric splitting's oddness and cubic coefficient
\eqref{eq:S2log} (through degree $12$), Campbell's conjugation series
\eqref{eq:campbell}, the Bernoulli inverse identity $\phi\beta=1$, and the
invariant-metric even-series identity \eqref{eq:metriceven} (through degree
$10$).
\end{itemize}
The supplied runs report $20$, $30$, and $26$ passed named checks
respectively; the script \texttt{code/run\_all\_verifiers.py} in the
article's directory reruns all three and records their reports. These are
exact finite algebraic verifications, not floating-point tolerance tests
and not a claim of formal proof-assistant certification. Finite-degree
agreement does not prove an all-orders identity, an analytic convergence
theorem, or a domain statement for an unbounded operator; those are
established by the proofs above.

\subsection{Two commutator encodings}
The right-nested data use exactly $R(w)$ of \eqref{eq:rightbracket}, so
that $Z_n=\sum_{|w|=n}\frac{c(w)}nR(w)$. This encoding is redundant: every
word ending in two equal letters represents zero, and the remaining brackets
are not claimed to be linearly independent.

For compact tables we use the standard Lyndon bracketing. Order the
alphabet by $X<Y$ and words lexicographically, a proper prefix preceding
its extensions. A nonempty word is \emph{Lyndon} if it is strictly smaller
than each of its nonempty proper suffixes. Define $L(X)=X$, $L(Y)=Y$, and
for a Lyndon word $w$ of length at least two write $w=uv$ with $v$ its
longest proper suffix that is Lyndon, and set
\begin{equation}\label{eq:lyndon}
 L(w)=[L(u),L(v)].
\end{equation}
(The standard factorization theorem guarantees that $u$ is again Lyndon;
the table generator applies the recursion to $u$ and $v$ in any case, so
every entry is an explicitly specified bracket polynomial.) For example
$L(\mathrm{XXY})=[X,[X,Y]]$ and $L(\mathrm{XYY})=[[X,Y],Y]$. Thus
$L(\mathrm{XYY})$ is \emph{not} $R(\mathrm{XYY})$, which is zero; confusing
the two conventions would corrupt many coefficients. No theorem about
Lyndon bases is needed to verify the tables: the generator expands each
tabulated bracket polynomial into words and requires the result to equal
the associative polynomial exactly. A table is therefore a certificate for
an explicit finite commutator identity whether or not one invokes the
general basis theorem. The counts in \cref{tab:counts} are numbers of
nonzero entries in these encodings, not dimensions of homogeneous free Lie
algebras and not counts of terms in a minimal basis.

\subsection{Data files and regeneration}
The principal machine-readable files in the article's \texttt{data/}
directory are listed in \cref{tab:files}. In an associative file a row
gives the coefficient of the literal word; in a \texttt{lyndon} file it
multiplies $L(w)$, and in the right-commutator file it multiplies $R(w)$.
Numerators and denominators are separate integers, and missing rows have
coefficient zero. The tree-polynomial file gives each $C_n$ as a polynomial
in the $A_j$ of \eqref{eq:An}; a key such as \texttt{2 3} denotes the
ordered product $A_2A_3$.
\begin{table}[htbp]
\centering\small
\begin{tabular}{@{}p{0.42\textwidth}p{0.52\textwidth}@{}}
\toprule
File & Meaning\\
\midrule
\texttt{bch\_associative.csv} & Every nonzero BCH word coefficient $c(w)$ through degree 12.\\
\texttt{bch\_right\_commutators.csv} & Dynkin-projection coefficients $c(w)/n$, with identically zero right brackets omitted.\\
\texttt{bch\_lyndon.csv} & Compact BCH commutator expansion through degree 12.\\
\texttt{zassenhaus\_associative.csv} & Associative coefficients of each $C_n$, $2\leq n\leq12$.\\
\texttt{zassenhaus\_lyndon.csv} & Compact commutator expansions of the same $C_n$.\\
\texttt{zassenhaus\_tree\_polynomials.json} & Each $C_n$, $2\leq n\leq10$, as a polynomial in the ordered $A_j$.\\
\texttt{strang\_lyndon.csv} & $\log(e^{X/2}e^Ye^{X/2})$ through degree 12.\\
\texttt{bernoulli\_plus.csv} & The coefficients $\bplus n=B_n^+/n!$ through $n=12$.\\
\bottomrule
\end{tabular}
\caption{Machine-readable companion expansions (exact rationals).}\label{tab:files}
\end{table}

The coefficient appendix (\cref{app:tables}) is generated from these
files by \texttt{code/make\_tables.py}, which independently recomputes every
printed degree-one to degree-six coefficient by the cut formula
\eqref{eq:wordcut} and by expansion of the displayed brackets, reconstructs
every printed Lyndon table from its brackets, and refuses to write the
appendix if any comparison fails. From the article's directory:
\begin{lstlisting}[language=bash]
python code/make_tables.py            # regenerate tex/appendix_tables.tex with checks
python code/run_all_verifiers.py      # rerun the three exact verifiers
latexmk -pdf -interaction=nonstopmode -halt-on-error bch_combined.tex
\end{lstlisting}
Python 3.10 or later with the standard library suffices; a standard TeX
installation with the packages named in the preamble compiles the source
without a bibliography processor, shell escape, or external files. The
verifiers accept larger cutoffs, but the number of words grows
exponentially with the degree; the finite coefficient formulas and
recurrences of this article, not the chosen degree limit, specify all
orders.
